\section{Primitive rows, calibrated prefixes, and the mixed packet}
\label{sec:interface-prefixes}

Fix \(h\in\Z\setminus\{0\}\), put \(H=\rad|h|\), and choose
\begin{equation}
 R=\lfloor X^{1/2-\delta}\rfloor,\qquad
 0<\delta<\frac12,\qquad R\ge H^2.
 \label{eq:R-def}
\end{equation}
In one separated primitive row packet,
\begin{equation}
 Q=LD,\qquad QJ\asymp X,\qquad
 L>\max\{2R,2|h|\},
 \label{eq:row-scales}
\end{equation}
and
\begin{equation}
 m=\ell d_0\asymp Q,\qquad
 \ell\in[L/2,2L]\text{ prime},\qquad
 d_0\in[D,2D],\qquad d_0\le R^{1/2}.
 \label{eq:row-support}
\end{equation}
The inherited primitive conditions include
\((d_0,H)=1\), \((j,H)=1\), and \(d_0H\le R\).  Hence
\begin{equation}
 (w,mH)=1
 \qquad\text{for every }w\mid mj+h.
 \label{eq:target-primitivity}
\end{equation}

The separated row weights satisfy
\begin{equation}
 \|x\|_\infty+\|y\|_\infty\ll_\eps X^\eps,\qquad
 \sum_m|x_m|^2+\sum_n|y_n|^2\ll QX^\eps,
 \label{eq:row-energy}
\end{equation}
and their absolute pair mass is denoted by
\begin{equation}
 \cZ_1:=\sum_{m,n}|x_my_n\mathfrak m(m,n)|
 \ll_\eps Q^2X^\eps.
 \label{eq:row-pair-mass}
\end{equation}
Here \(\mathfrak m\) is the actual bounded generic row-pair mask.  We
use only
\begin{equation}
 |\mathfrak m(m,n)|\le1,\qquad
 \mathfrak m(m,m)=0,\qquad
 \mathfrak m\text{ is independent of all divisor variables}.
 \label{eq:mask}
\end{equation}
The last two properties are essential in different places: the
conductor tensor uses divisor independence, while the large-\(g\)
zero-mode argument uses \(m-n\ne0\).

For the strongest logarithmic conclusion we state separately the
fixed-projective-mass hypothesis
\begin{equation}
 \cZ_1\ll Q^2(\log X)^K
 \label{eq:polylog-projective-mass}
\end{equation}
for some fixed \(K\ge0\).  If a scalar packet is reassembled from a
projective sum \(\sum_\nu x^{(\nu)}\otimes y^{(\nu)}\), the left side
of \eqref{eq:polylog-projective-mass} is replaced by
\(\sum_\nu\|x^{(\nu)}\|_1\|y^{(\nu)}\|_1\).  The soft estimate
\eqref{eq:row-pair-mass} does not imply
\eqref{eq:polylog-projective-mass}; the two hypotheses are kept
distinct throughout.

\subsection{Fourier and fiber notation}

We use
\begin{equation}
 e(t)=e^{2\pi it},\qquad e_q(t)=e(t/q),\qquad
 c_q(r)=\sum_{\xi\in(\Z/q\Z)^\times}e_q(r\xi),
 \label{eq:additive-normalizations}
\end{equation}
and
\begin{equation}
 \widehat F(t)=\int_{\R}F(x)e(-tx)\,dx.
 \label{eq:Fourier-normalization}
\end{equation}
Write \(\cU_q=(\Z/q\Z)^\times\) and let
\(\mathfrak X(q)\) be its group of Dirichlet characters.  Whenever
\((H,q)=1\), the symbol \(\overline H\) denotes the inverse of \(H\)
modulo \(q\).  Finally, the centered divisor indicator is
\begin{equation}
 C_{m,w}(j)=\one_{w\mid mj+h}-\frac1w.
 \label{eq:centered-divisor-indicator}
\end{equation}

\subsection{Divisor coefficients and exact densities}

Put
\begin{equation}
 a(w)=-\mu(w)\log w.
 \label{eq:a-def}
\end{equation}
The finite Ramanujan model has the divisor expansion
\begin{equation}
 \Lambda_R(n)=\sum_{w\mid n}\lambda'_R(w),
 \qquad
 \lambda'_R(w)=
 w\sum_{b\le R/w}\frac{\mu(wb)\mu(b)}{\varphi(wb)}
 \quad(w\le R),
 \label{eq:lambda-def}
\end{equation}
and \(\lambda'_R(w)=0\) for \(w>R\).  Define
\begin{equation}
 b_R(w)=a(w)-\lambda'_R(w).
 \label{eq:b-def}
\end{equation}
All three coefficients are supported on squarefree integers, and
uniformly for \(w\le X^{O(1)}\),
\begin{equation}
 |a(w)|+|b_R(w)|+|\lambda'_R(w)|
 \ll_\eps X^\eps.
 \label{eq:coefficient-soft}
\end{equation}

For a row \(m=\ell d_0\), the exact finite-model and calibrated prime
densities are
\begin{equation}
 \rho_{M,H,R}(m)=
 \frac{H}{\varphi(H)}\frac {d_0}{\varphi(d_0)},
 \qquad
 \rho_{P,H}(m)=\frac{mH}{\varphi(mH)}.
 \label{eq:two-densities}
\end{equation}
TPC-19--20 prove the finite identity
\begin{equation}
 \sum_{\substack{w\le R\\(w,mH)=1}}
 \frac{\lambda'_R(w)}w
 =\rho_{M,H,R}(m)
 \label{eq:model-density-identity}
\end{equation}
and the row drift formula
\begin{equation}
 \delta_H(m):=\rho_{P,H}(m)-\rho_{M,H,R}(m)
 =
 \frac{H}{\varphi(H)}
 \frac{d_0}{\varphi(d_0)(\ell-1)}
 \ll_{\eps,h}\frac{X^\eps}{L}.
 \label{eq:row-drift}
\end{equation}
These are exact finite algebra, not prime distribution in a growing
progression \citep{WangTPC19,WangTPC20}.

\subsection{Uniform endpoint calibration}

For \(M\ge1\) and \(Y\ge1\), set
\begin{equation}
 A_M(Y)=
 \sum_{\substack{w\le Y\\(w,M)=1}}
 \frac{-\mu(w)\log w}{w},
 \qquad
 B_M(Y)=
 \sum_{\substack{w\le Y\\(w,M)=1}}\frac{\mu(w)}w.
 \label{eq:ABM}
\end{equation}

\begin{lemma}[Uniform calibration and its companion]
\label{lem:uniform-companion}
For every fixed \(A,C>0\), uniformly for \(Y\ge3\) and
\(M\le Y^C\),
\begin{align}
 A_M(Y)
 &=
 \frac M{\varphi(M)}
 +O_{A,C}\left(
 \frac M{\varphi(M)}(\log Y)^{-A}\right),
 \label{eq:A-calibration}\\
 B_M(Y)
 &\ll_{A,C}
 \frac M{\varphi(M)}(\log Y)^{-A}.
 \label{eq:B-calibration}
\end{align}
Consequently, if \(Y_2\ge Y_1\ge X^\kappa\) and
\(M\le X^{C_0}\), where \(C_0,\kappa>0\) are fixed, then for every
\(A>0\),
\begin{align}
 A_M(Y_2)-A_M(Y_1)
 &\ll_{A,C_0,\kappa}(\log X)^{-A+1},
 \label{eq:A-interval}\\
 B_M(Y_2)-B_M(Y_1)
 &\ll_{A,C_0,\kappa}(\log X)^{-A+1},
 \label{eq:B-interval}
\end{align}
uniformly with constants depending on \(A,C_0,\kappa\).
\end{lemma}

\begin{proof}
Equation \eqref{eq:A-calibration} is the uniform M\"obius calibration
theorem of TPC-20 \citep{WangTPC20}.  We recall why the same proof gives
\eqref{eq:B-calibration}.  Let
\[
 \mathscr S(M)=\{k\ge1:p\mid k\Longrightarrow p\mid M\},
 \qquad
 G_M(s)=\prod_{p\mid M}(1-p^{-s})^{-1}.
\]
The exact convolution
\[
 \mu(n)\one_{(n,M)=1}=(\mu*\one_{\mathscr S(M)})(n)
\]
gives
\begin{equation}
 B_M(Y)=
 \sum_{\substack{k\le Y\\k\in\mathscr S(M)}}
 \frac1k
 \sum_{r\le Y/k}\frac{\mu(r)}r.
 \label{eq:B-convolution}
\end{equation}
The classical zero-free region for \(\zeta(s)\) implies
\[
 \sum_{r\le Z}\frac{\mu(r)}r\ll_A(\log Z)^{-A}
 \qquad(Z\ge3).
\]
Split \eqref{eq:B-convolution} at
\(k=Y/\exp(\sqrt{\log Y})\).  The central range is bounded with the
displayed zero-free-region estimate.  The remaining range is handled
by the same Rankin bound for \(G_M(1-\eta)\) used in the proof of
\eqref{eq:A-calibration}; it is uniform for \(M\le Y^C\).
This proves \eqref{eq:B-calibration}.

Apply the first part with \(C=C_0/\kappa\), and subtract
\eqref{eq:A-calibration} at \(Y_2\) and \(Y_1\): the common
main term \(M/\varphi(M)\) cancels.  Do the same with
\eqref{eq:B-calibration}, use
\(M/\varphi(M)\ll\log\log(3M)\), and increase \(A\) to absorb this
factor.  Since \(\log Y_1\ge\kappa\log X\), the interval estimates
follow.  The only external analytic input is the classical zero-free
region; see \citet[Chapter~5]{IwaniecKowalski2004}.
\end{proof}

\subsection{Complete prefixes}

For \(R\le W\le U\), define
\begin{equation}
 \vartheta_W(w)=b_R(w)\one_{w\le W},
 \qquad
 E_W(m)=A_{mH}(W)-\rho_{P,H}(m),
 \label{eq:theta-E}
\end{equation}
and
\begin{equation}
 P_{m,W}(j)=
 E_W(m)+
 \sum_{\substack{w\le W\\(w,mH)=1}}
 b_R(w)C_{m,w}(j).
 \label{eq:calibrated-prefix}
\end{equation}

\begin{proposition}[Exact calibrated-prefix identities]
\label{prop:prefix-identities}
For every \(R\le W\le U\),
\begin{align}
 \sum_{\substack{w\le W\\(w,mH)=1}}\frac{b_R(w)}w
 &=\delta_H(m)+E_W(m),
 \label{eq:prefix-density}\\
 P_{m,W}(j)
 &=
 \sum_{\substack{w\le W\\w\mid mj+h}}b_R(w)-\delta_H(m).
 \label{eq:prefix-indicator}\\
 P_{m,T}(j)-P_{m,S}(j)
 &=
 \sum_{\substack{S<w\le T\\w\mid mj+h}}a(w).
 \label{eq:prefix-increment}
\end{align}
If \(U\) contains every positive target in the packet, then
\begin{equation}
 P_{m,U}(j)=
 \Lambda(mj+h)-\Lambda_R(mj+h)-\delta_H(m),
 \label{eq:full-residual-prefix}
\end{equation}
which is the locally centered residual signal of TPC-20.
\end{proposition}

\begin{proof}
Since \(\lambda'_R\) is supported on \(w\le R\),
\[
 \sum_{\substack{w\le W\\(w,mH)=1}}\frac{b_R(w)}w
 =A_{mH}(W)-\rho_{M,H,R}(m).
\]
Insert \eqref{eq:two-densities} to obtain
\eqref{eq:prefix-density}.  Expanding
\(C_{m,w}=\one_{w\mid mj+h}-1/w\) in
\eqref{eq:calibrated-prefix} then proves
\eqref{eq:prefix-indicator}.  Subtract the formulas at \(T\) and
\(S\); on \(w>R\), \(b_R(w)=a(w)\), proving
\eqref{eq:prefix-increment}.  At the target horizon, the von Mangoldt
divisor identity and \eqref{eq:lambda-def} give
\eqref{eq:full-residual-prefix}.
\end{proof}

\begin{lemma}[Prefix coefficient mass]
\label{lem:prefix-mass}
For every fixed \(C>0\), uniformly for \(R\le W\le X^C\),
\begin{equation}
 \sum_{w\le W}\frac{|b_R(w)|}{w}\ll_C(\log(2X))^3.
 \label{eq:prefix-L1}
\end{equation}
Moreover, uniformly in squarefree \(g\le W\),
\begin{equation}
 \sum_{\substack{c\le W/g\\(c,g)=1}}
 \frac{|b_R(gc)|}{c}\ll_C(\log(2X))^4.
 \label{eq:cofactor-L1}
\end{equation}
\end{lemma}

\begin{proof}
The \(a\)-part is bounded by a harmonic logarithmic sum.  For the
finite-model part use the positive formula
\[
 |\lambda'_R(w)|
 \le\frac w{\varphi(w)}
 \sum_{b\le R/w}\frac{\mu(b)^2}{\varphi(b)}
 \qquad(w\text{ squarefree})
\]
and sum the resulting nonnegative double divisor series.  This gives
\(\sum_{w\le R}|\lambda'_R(w)|/w\ll(\log(2R))^2\), hence
\eqref{eq:prefix-L1}.  In \eqref{eq:cofactor-L1}, use
\(\varphi(gc)=\varphi(g)\varphi(c)\) on \((g,c)=1\), and bound the
two remaining harmonic divisor sums.  The harmless factor
\(g/\varphi(g)\) is absorbed by one additional logarithm.
\end{proof}

\subsection{The calibrated mixed increment and inherited interfaces}

Let \(F\) range over a uniformly bounded subset of
\(C_c^\infty(\R)\), and put
\begin{equation}
 L_{m;S,T}(j):=P_{m,T}(j)-P_{m,S}(j)
 =\sum_{\substack{S<w\le T\\w\mid mj+h}}a(w).
 \label{eq:long-increment}
\end{equation}
The main packet is
\begin{align}
 \cI_{S,T}
 :={}&
 \sum_{m,n}x_my_n\mathfrak m(m,n)
 \sum_{\substack{j\in\Z\\(j,H)=1}}F(j/J)
 P_{m,S}(j)L_{n;S,T}(j).
 \label{eq:mixed-packet}
\end{align}
Its natural scale is
\begin{equation}
 Q^2J\asymp XQ.
 \label{eq:natural-scale}
\end{equation}

For a compatible pair \(u,v\), let \(q=[u,v]\) and let
\(\rho_{m,n}^{u,v}\in\cU_q\) be the joint residue class determined by
\[
 u\mid m\rho_{m,n}^{u,v}+h,\qquad
 v\mid n\rho_{m,n}^{u,v}+h.
\]
Define the exact joint fiber
\begin{align}
 Z_{q;S,T}(\xi)
 :={}&
 \sum_{\substack{u\le S,\ S<v\le T\\{}[u,v]=q}}
 b_R(u)a(v)
 \sum_{\substack{m,n\\(u,v)\mid m-n\\
                  \rho_{m,n}^{u,v}\equiv\xi\pmod q}}
 x_my_n\mathfrak m(m,n),
 \label{eq:joint-fiber}
\end{align}
with all row and target-primitivity restrictions retained.  Put
\begin{equation}
 \overline Z_{q;S,T}
 =\frac1{\varphi(q)}\sum_{\xi\in\cU_q}Z_{q;S,T}(\xi),
 \qquad
 D_{q;S,T}=Z_{q;S,T}-\overline Z_{q;S,T}.
 \label{eq:joint-fiber-centering}
\end{equation}
The joint nonzero scalar packet is
\begin{align}
 \mathscr J_{q;S,T}^{\ne0}
 ={}&\frac{J}{Hq}\sum_{r\ne0}c_H(r)
 \widehat F\left(\frac{rJ}{Hq}\right)
 \sum_{\xi\in\cU_q}Z_{q;S,T}(\xi)
 e_q(r\overline H\xi).
 \label{eq:joint-nonzero-scalar}
\end{align}
Its mean and discrepancy parts are obtained by replacing \(Z_{q;S,T}\)
with \(\overline Z_{q;S,T}\) and \(D_{q;S,T}\), respectively.

We assume explicitly the TPC-20 no-wrap condition
\begin{equation}
 J\ge4HX^{\epsilon_0}
 \label{eq:no-wrap}
\end{equation}
for some fixed \(\epsilon_0>0\) \citep{WangTPC20}.  Fix once and for
all
\begin{equation}
 0<\omega<\epsilon_0.
 \label{eq:Fourier-window-margin}
\end{equation}
To state the Fourier-tail interface quantitatively, define
\begin{align}
 \mathscr T_\omega
 :={}&
 \sum_{q\le ST}\frac{J}{Hq}
 \sum_{|r|>X^\omega Hq/J}|c_H(r)|
 \left|\widehat F\left(\frac{rJ}{Hq}\right)\right|
 \sum_{\xi\in\cU_q}|Z_{q;S,T}(\xi)|.
 \label{eq:Fourier-tail-mass}
\end{align}
The inherited polynomial tail-mass hypothesis is that for every fixed
\(A>0\), the Fourier decay order may be chosen so that
\begin{equation}
 \mathscr T_\omega\ll_{A,F,H}X^{-A}\cZ_1.
 \label{eq:Fourier-tail-interface}
\end{equation}
This is the precise tail discarded in the cell reassembly below.  The
fixed inequality \(\omega<\epsilon_0\), not an \(A\)-dependent window,
is what preserves no-wrap on the retained frequencies.
Moreover,
\[
 \sum_{\xi\in\cU_q}|D_{q;S,T}(\xi)|
 \le 2\sum_{\xi\in\cU_q}|Z_{q;S,T}(\xi)|,
\]
so \eqref{eq:Fourier-tail-interface} controls the discrepancy tail as
well as the undecomposed joint tail.  On the retained window,
\(|r|\le X^\omega Hq/J<q/4\) for all sufficiently large \(X\), by
\eqref{eq:no-wrap} and \eqref{eq:Fourier-window-margin}.

For the four-channel interpretation in
\cref{sec:asymmetric-four-channel} we also record the short-side margin
\begin{equation}
 HS\le JX^{-\eta_0}
 \label{eq:short-side}
\end{equation}
for some fixed \(\eta_0>0\).  The exact calibrated recombination and
the conductor theorem themselves do not require
\eqref{eq:short-side}.  No comparison between \(T\) and \(J\), or
between \(L\) and \(T\), is imposed.

\begin{proposition}[Exact joint-minus-drift recombination]
\label{prop:joint-minus-drift}
Pointwise in the two rows and the orbit variable,
\begin{align}
 P_{m,S}(j)L_{n;S,T}(j)
 ={}&
 \sum_{\substack{u\le S\\u\mid mj+h}}b_R(u)
 \sum_{\substack{S<v\le T\\v\mid nj+h}}a(v)
 -\delta_H(m)L_{n;S,T}(j).
 \label{eq:joint-minus-drift}
\end{align}
Moreover, for every fixed \(\epsilon>0\), the drift contribution to
\eqref{eq:mixed-packet} is
\begin{equation}
 \ll_{\epsilon,F,h}J\cZ_1L^{-1}X^\epsilon.
 \label{eq:drift-packet-bound}
\end{equation}
\end{proposition}

\begin{proof}
The identity follows at once from
\eqref{eq:prefix-indicator} and \eqref{eq:long-increment}.  Also
\[
 |L_{n;S,T}(j)|
 \le \sum_{v\mid nj+h}|a(v)|
 \le \tau(nj+h)\log(2X)\ll_{\epsilon,F}X^\epsilon.
\]
There are \(O_F(J)\) orbit points, and
\eqref{eq:row-drift} and the definition of \(\cZ_1\) prove
\eqref{eq:drift-packet-bound} after renaming \(\epsilon\).
\end{proof}
